%% New owned fragment; released source01 manuscript files remain frozen.
\subsubsection{The exact selected-mono localization and its logical functor}
\label{sec:p10-lawvere-localization}

The original article's datum is the finite-limit-preserving
\(u_*:\mathcal E\to\mathcal E_0\) and the internal Heyting
homomorphism \(h:u_*(T)\to T_0\), not an assumption of a
geometric inverse image or an ordinary model-theoretic ultrafilter
\cite[p.~332]{Lawvere1970}.  The following construction supplies
the ordinary localization and its logical preservation at exactly
that datum.  Its reduction to a truth filter is proved by the
universal property; no finite, principal, Boolean or prime-filter
restriction is imposed.  The base-relative closed-category and
collapse assertions remain distinct below.

\begin{hbdefinition}{def:p10-lawvere-localization-filter}{The full truth filter induced by the original datum}
Identify \(\operatorname{Sub}_{\mathcal E}(1)\) with the
maps \(1\to T\) by classification.  Define
\[
 \theta_h(U)=h\circ u_*(\chi_U):1_0\longrightarrow T_0,
 \qquad
 \mathcal F_h=\{U\hookrightarrow1:\theta_h(U)=\mathsf{true}_0\},
\]
using the finite-limit identification \(u_*(1)\cong1_0\).
The value \(\theta_h(U)\) is not asserted to be the
characteristic map of \(u_*(U)\).  For a mono
\(m:A\hookrightarrow X\), put
\[
         \gamma_m=\forall_{X\to1}(m)\in\operatorname{Sub}(1),
 \qquad \mathcal W_h=\{m:\gamma_m\in\mathcal F_h\}.
\]
The predicate of \(\gamma_m\) is exactly
\(\sigma_1(\Pi_{X\to1}(m))\), by the proved restriction
of the family adjunction to monos.  Thus \(\mathcal W_h\)
is precisely the monomorphism class in the source display,
without interchanging \(u_*\) with universal quantification.
\end{hbdefinition}

\begin{hblemma}{thm:p10-lawvere-localization-denominators}{proved=annex:p10.lawvere-localization.denominators}{The source denominators have uniform cofinal restrictions}
The map \(\theta_h\) is a Heyting homomorphism on global
truth algebras.  The set \(\mathcal F_h\) contains top,
is upward closed, and is closed under binary intersection.
It is proper when \(\mathcal E_0\) is nondegenerate;
the degenerate-base case permits the improper filter.
For every \(U\in\mathcal F_h\), the mono
\(i_{U,X}:U\times X\hookrightarrow X\) lies in
\(\mathcal W_h\).  Conversely every \(m\in\mathcal W_h\)
contains \(\gamma_m\times X\), and its pullback to
that admitted subterminal is an isomorphism.
Hence localization at the source's full mono class and
localization at all \(i_{U,X}\), \(U\in\mathcal F_h\),
have the same universal property.  This is a mathematical
equivalence of the exact classes, not replacement of the
general datum by a special finite filter.
\end{hblemma}

\begin{hbdefinition}{def:p10-lawvere-localization-fractions}{Fractions and their actual equality relation}
Fix the locally small/universe convention, so \(\operatorname{Sub}(1)\)
and each indicated hom collection are sets.  The objects of
\(\mathcal E/\mathcal F_h\) are those of \(\mathcal E\).
A morphism \(X\to Y\) is represented by
\((U,f)\), where \(U\in\mathcal F_h\) and
\(f:U\times X\to Y\).  Set
\[
 (U,f)\sim(V,g)\quad\Longleftrightarrow\quad
 \text{some }W\in\mathcal F_h, W\leq U\cap V,
       \text{ has }f|_{W\times X}=g|_{W\times X}.
\]
Write \([U,f]\) for its class.  Identity has denominator
\(1\).  Compose \([U,f]:X\to Y\) and
\([V,g]:Y\to Z\) on \(W=U\cap V\): restrict
\(f\) to \(W\times X\), pair it with the unique
parameter map to \(V\), then apply \(g\).
No representatives outside the admitted subterminal are
chosen or required.  The canonical functor is
\(Q(X)=X\), \(Q(f)=[1,f]\).

In the 1992 primary comparison, a representative is a
map \(X\times U\to Y\times U\) over \(U\).
Its first coordinate is the present representative;
the other coordinate is forced because \(U\) is
subterminal.  Product symmetry gives the exact typed
bridge.  The ordinary set-sized quotient of hom sets
does not assume an infinite colimit inside \(\mathcal E\)
\cite[V\S9, pp.~258--260]{MacLaneMoerdijk1992}.
That source's filter convention excludes zero; the
improper-filter boundary here is explicitly separate.
\end{hbdefinition}

\begin{hbtheorem}{thm:p10-lawvere-localization-fractions}{proved=annex:p10.lawvere-localization.fractions}{The complete fraction calculus and localization property}
The displayed relation is an equivalence relation and
the composition is independent of representatives and
refinements, associative and unital.  The functor \(Q\)
inverts every source mono in \(\mathcal W_h\).
For any category \(\mathcal D\) and functor
\(H:\mathcal E\to\mathcal D\) inverting that class,
there is a unique extension on the same objects,
\[
 \bar H[U,f]=H(f)\,H(i_{U,X})^{-1}.
\]
Natural transformations extend uniquely by their original
components.  Consequently the constructed category is
the ordinary \(\mathcal E_h\) formally localized at
the exact class stated on p.~332.  Every source fraction
has a uniform representative of this form, and equality
is exactly agreement on a further admitted common
restriction; there is no additional unproved fraction
existence or coherence requirement.
\end{hbtheorem}

\begin{hbtheorem}{thm:p10-lawvere-localization-logical}{proved=annex:p10.lawvere-localization.logical}{The ordinary localized topos and full logical preservation}
The constructed \(\mathcal E_h\) has finite limits,
finite colimits, exponentials and subobject classifier
\(Q(\mathsf{true}):Q1\to QT\), and \(Q\) preserves
all these structures up to their canonical isomorphisms.
In particular \(Q(Y^X)\) is the exponential of
\(QY\) by \(QX\), with its actual evaluation, and
every localized mono has a unique characteristic map,
not merely those arising from initially chosen strict
monos.  Every finite diagram of germs is handled on
a common admitted restriction, including equalizers
and coequalizers of genuinely localized arrows.

Power objects, membership, the graph constructions of
dependent products, and both internal quantifier maps
are preserved as well.  Thus \(Q\) is a logical
morphism in the original source's package.  The
improper filter gives the degenerate topos and is
included explicitly.  The later primary comparison is
\cite[V\S9, Theorem~2 and Lemma~3, pp.~261--263]
{MacLaneMoerdijk1992}; its proof boundaries and proper
filter convention are not silently promoted to a
different 1970 source statement.
\end{hbtheorem}

\begin{hbexample}{def:p10-lawvere-localization-base-functor}{The original base functor need not descend unchanged}
Let \(\mathcal E=\mathsf{Set}\times\mathsf{Set}\),
\(\mathcal E_0=\mathsf{Set}\), \(u_*\) be binary
product, and \(h:2\times2\to2\) project to the first
coordinate.  These satisfy the original hypotheses.
The subterminal \(U=(1,0)\) is admitted, so its mono
into \(1=(1,1)\) is inverted.  But its image under
\(u_*\) is \(0\to1\), not an isomorphism.
Therefore \(u_*\) itself cannot factor through this
localization.  A new base-relative construction is
needed; finite-limit preservation alone does not
produce a geometric direct image or a strong
closed-functor descent.
\end{hbexample}

\begin{proofboundary}{def:p10-lawvere-localization-boundary}{The exact completed and still-owed source assertions}
The ordinary localization and its logical topos
preservation are proved at the full original \(u_* / h\)
datum, through an exact cofinal reduction to the
arbitrary induced filter.  The original article only
announces these conclusions.  Its additional claim
that \(\mathcal E_h\) is defined over \(\mathcal E_0\)
in the sense of closed categories still requires its
actual base-relative construction and precise closed
structure; it is not supplied by asserting that
\(u_*\) descends.  The two-valued collapse and broader
BVM/classification/CH dependencies remain separate
full-proof obligations.  No improper/principal/finite
special case closes those general obligations.
\end{proofboundary}
